\cotefile{161-2}
\scanpage{51}
\inserted{(\unsure{?? sens|22 suit})}
\deleted{\unsure{S}}Soit $\Pi'_n=\bigcap\Pi_{nj}$, et $e'_n=\operatorname{Im}(\Pi'_n\mathbin{\unsure{-|\to}}e_E)$. La condition est que $e'_n\to e_E$ est un \uline{iso}, i.e. $\pi'_n\to 1_E$ un épi. Il n'est pas clair cependant que la formation des $e'_n$ commute aux chgt de topos de base $E'\mathbin{-}E$, car ce n'est pas clair pour la formation des $\Pi_{nj}$. Donc il n'est pas clair qu'on a un topos modulaire \deleted{\unsure{triplement}},

Mais prenons d'abord le \unsure{pbt|pbl|pts} analogue des anneaux qui sont \uline{de plus} de support tout.
\margin{\uline{lemme} \uline{à vérifier}}% slanted
Alors on voit que localement le choix des $e_i$ est bien déterminé, \deleted{\unsure{à cela près|à cette fois}} à permutation près, et de plus que si \deleted{\unsure{on a un}} on désigne localement par $m$ le nb de $e_i$ qui sont $\ne 0$, alors \deleted{la décomposition} les $e_i$ sont tous \deleted{\unsure{à}} \unsure{différent|différents} (\unsure{et|c'est} pour cause, $e_i=e_j$ \unsure{et} $e_i^2=e_i$, $e_i e_j(=e_i^2)=0$ implique $e_i=0$!) Si on se borne pour le moment au cas où $n=m$, on trouve que \deleted{le choix des} \deleted{\unsure{le}} faisceau des choix possibles des \deleted{$\unsure{e}$} $e_i$ est un pseudo-torseur \inserted{sur $E$} sous le groupe \inserted{symétrique} $\unsure{\mathfrak S|\mathscr S}_n$; dire \deleted{que} \deleted{est} $A$ est \inserted{quasi} connexe de rang \uline{exact} $n$, (ce qui implique, par définition, à support \deleted{\unsure{et $\ne$}} tout). \deleted{\unsure{Ceci}}

Je ne sais si les \uline{pseudo}-\uline{torseurs} commutent à tt chgt de base \inserted{$\unsure{E'|E_i'}\mathbin{-}E$}, mais le torseur oui, si on se borne aux $E$ munis de $A$ quasi-connexe de rang exact $n$. Or la donnée d'un tel revient à la donnée

a) d'un $\mathfrak S_{n\,E}$-torseur $P$ (sur $E$, i.e. d'un morphisme $E\xrightarrow{p}B_{\mathfrak S_n}$ (\unsure{ou plus de|ou plus des} $E\to\widehat R=B)$,

b) d'une décomposition de $A_P$ en produit $\prod_{1\le i\le n} A_i$ avec les $A_i$ connexes (néc. à support tt), et avec compatibilité entre celle-ci et l'action de $\mathfrak S_n$.

Cela signifie qu'on a le torseur $P$ (i.e. $p$), des anneaux connexes \deleted{\unsure{non}} à support tout $A_i$ situés sur $P$, et \deleted{\unsure{une}} action de $\mathfrak S_n$ sur $P$ permutant \unsure{entre eux|avec} les $A_i$ suivant des actions sur les indices.
% archive sheet number: 50
